\section{Bibliografía}
% Formato APA 7. Todas las entradas tienen que estar citadas en el cuerpo del texto.
% La lista usa sangría francesa de 1.27 cm, como pide APA 7 (§2.12).
\begin{list}{}{%
  \setlength{\leftmargin}{1.27cm}%
  \setlength{\itemindent}{-1.27cm}%
  \setlength{\labelwidth}{0pt}%
  \setlength{\labelsep}{0pt}%
  \setlength{\itemsep}{0.6em}%
  \setlength{\parsep}{0pt}%
}
  \item[] Anthropic. (2026). \textit{Claude Code} (Opus 5) [Modelo de lenguaje grande]. Recuperado el 9 de septiembre de 2026, de \url{https://claude.ai/code}

  \item[] Goldberg, D. (1991). What every computer scientist should know about floating-point arithmetic. \textit{ACM Computing Surveys, 23}(1), 5--48. \url{https://doi.org/10.1145/103162.103163}

  \item[] Harris, C. R., Millman, K. J., van der Walt, S. J., Gommers, R., Virtanen, P., Cournapeau, D., Wieser, E., Taylor, J., Berg, S., Smith, N. J., Kern, R., Picus, M., Hoyer, S., van Kerkwijk, M. H., Brett, M., Haldane, A., del Río, J. F., Wiebe, M., Peterson, P., \ldots{} Oliphant, T. E. (2020). Array programming with NumPy. \textit{Nature, 585}(7825), 357--362. \url{https://doi.org/10.1038/s41586-020-2649-2}

  \item[] Higham, N. J. (2002). \textit{Accuracy and stability of numerical algorithms} (2.\textsuperscript{a} ed.). Society for Industrial and Applied Mathematics. \url{https://doi.org/10.1137/1.9780898718027}

  \item[] Hunter, J. D. (2007). Matplotlib: A 2D graphics environment. \textit{Computing in Science \& Engineering, 9}(3), 90--95. \url{https://doi.org/10.1109/MCSE.2007.55}

  \item[] Institute of Electrical and Electronics Engineers. (2019). \textit{IEEE standard for floating-point arithmetic} (IEEE Std 754-2019). \url{https://doi.org/10.1109/IEEESTD.2019.8766229}

  \item[] International Organization for Standardization. (2018). \textit{Information technology. Programming languages. C} (ISO/IEC 9899:2018).

  \item[] Kahan, W. (1965). Pracniques: Further remarks on reducing truncation errors. \textit{Communications of the ACM, 8}(1), 40. \url{https://doi.org/10.1145/363707.363723}

  \item[] Python Software Foundation. (2024). \textit{The Python language reference} (Versión 3.12). Recuperado el 9 de septiembre de 2026, de \url{https://docs.python.org/3.12/reference/}

  \item[] Sterbenz, P. H. (1974). \textit{Floating-point computation}. Prentice-Hall.
\end{list}
